\documentclass[a4paper,11pt]{article}

% - Encoding and language -
\usepackage[T2A]{fontenc}
\usepackage[utf8]{inputenc}
\usepackage[russian]{babel}

% - Page layout and formatting -
\usepackage[left=1.2cm,right=1.2cm,top=1.2cm,bottom=1.2cm]{geometry}
\usepackage{titlesec}
\usepackage{enumitem}
\usepackage{hyperref}
\usepackage{xcolor}
\usepackage{parskip}

% - Remove paragraph indentation -
\setlength{\parindent}{0pt}
\setlength{\parskip}{2pt}

% - Link colors -
\hypersetup{
  colorlinks=true,
  urlcolor=blue,
  linkcolor=black
}

% - Section formatting -
\titleformat{\section}
  {\large\bfseries}
  {}
  {0pt}
  {}
  [\titlerule]

\titlespacing{\section}{0pt}{0.6em}{0.3em}
\setlist[itemize]{leftmargin=1.2cm, itemsep=1pt}

\begin{document}
\pagestyle{empty}

% - Header -
\begin{center}
{\huge \textbf{Даниил Спиридонов}}\\[4pt]
Machine Learning Engineer $\cdot$ Москва, Россия\\[2pt]
Телефон: +7 (936) 505-12-45 | Telegram: \href{https://t.me/daspiridonov}{@daspiridonov} | \href{mailto:sprv-4@yandex.ru}{sprv-4@yandex.ru}\\
\href{https://www.linkedin.com/in/daniel-spiridonov/}{linkedin.com/in/daniel-spiridonov} | \href{https://github.com/notaskynet}{github.com/notaskynet}
\end{center}

% - About -
\section*{О себе}
ML-инженер, 1,5 года коммерческого опыта. Основные направления: LLM-сервисы и агентские пайплайны, классический ML, MLOps. Студент 5 курса НИУ ВШЭ.

% - Experience -
\section*{Опыт работы}
\textbf{РТ-ИБ -- Инженер машинного обучения}\hfill \textit{Декабрь 2024 -- н.\,в.}

\begin{itemize}
  \item Разработал и вывел в прод LLM-агента для triage инцидентов на LangGraph, сократив время разбора инцидента аналитиком: агент готовит структурированный разбор перед эскалацией
  \item Автоматизировал написание правил-исключений через LLM-пайплайн на self-hosted vLLM, покрыв \(\sim\)70\% случаев и заменив ручную работу аналитика на ревью сгенерированных правил
  \item Обучил и вывел в прод классификатор алертов на CatBoost, отсеивающий \(\sim\)30\% ложных срабатываний при сохранении 99\% истинных
  \item Внедрил детектор аномалий по \(\sim\)100 источникам событий с автовыбором метода, работающий в проде с интервалом 10 минут
  \item Переписал сервис группировки инцидентов, покрывающий \(\sim\)15\% всех группировок, ускорив замену и сравнение ретривал-моделей
  \item Выстроил MLOps-инфраструктуру команды (CI/CD, версионирование датасетов и артефактов, стандарты code quality), ускорив старт новых ML-проектов
  \item Проводил технические собеседования кандидатов и онбординг новых членов команды
\end{itemize}

\textbf{GameGlory (\href{https://gameglory.studio/}{gameglory.studio}) -- ML-инженер, проектная работа}\hfill \textit{Октябрь 2025 -- н.\,в.}

\begin{itemize}
  \item Разработал семантический поиск конкурентов по описанию и тегам игры
  \item Разработал LLM-анализ идеи: скоринг по ключевым аспектам с рекомендациями и оценкой рисков
  \item Построил сбор данных по каталогу игр Steam и оценку недоступных напрямую метрик (выручка, класс издателя)
\end{itemize}

% - Skills -
\section*{Технические навыки}
\begin{itemize}
  \item \textbf{Языки программирования:} Python, SQL
  \item \textbf{Фреймворки и библиотеки:} pandas, numpy, scikit-learn, CatBoost, XGBoost, PyTorch, FastAPI, LangGraph, vLLM, Langfuse
  \item \textbf{MLOps и инфраструктура:} Docker, GitLab CI/CD, Airflow, MLflow, MinIO, PostgreSQL, RabbitMQ, Grafana
  \item \textbf{Dev-инструменты:} git, ruff, mypy, uv
  \item \textbf{Иностранные языки:} Английский (B2)
\end{itemize}

% - Pet projects -
\section*{Pet-проекты}
\begin{itemize}
  \item \href{https://github.com/notaskynet/snowcite}{\textbf{snowcite}} -- MCP-сервер для литературных обзоров: поиск по arXiv/Semantic Scholar/OpenAlex/PubMed, сборка PDF через LaTeX/Typst
\end{itemize}

% - Education -
\section*{Образование}
\textbf{НИУ ВШЭ} -- специалитет «Компьютерная безопасность»\\
\textit{2022 -- 2028 (ожидается)}

\end{document}
